\section{Introduction}
Informally we describe Linear Representation Hypothesis (LRH) to be that concepts (e.g Gender, Vehicle, etc)
can be represented by linear directions in the activation space, \citep{Mikolov2013LinguisticRI}. 
We may also define LRH formally, following \citet{park2023LRH}. We start by defining concepts, 
which we say are variables that can be changed in isolation.
And then LRH says that each concept corresponds to a particular linear subspace (Or direction).
Throughout this project we will switch between both as convenient and depending on the context,
gathering empirical evidence to also evaluate if our findings are consistent with the theoretical foundations
of LRH as presented in \citet{park2023LRH}. Furthermore, Large Language Models (LLMs) are susceptible 
to malicious prompts and may produce responses based on harmful stereotypical associations
\citet{jeoung-etal-2023-stereomap}. While there are many safety measurements in place to combat this, 
they tend to come at the cost of the model's own utility (i.e overtuning for safety and harming utility).
Safety and utility both satisfy \citet{park2023LRH} defintions of concepts, so we may reasonablly ask if they are
causally separable.
We use both notions of LRH, and both suggest that safety and utility may be represented by linear directions or for 
a theoretical basis, they may occupy orthogonal subspaces. So we aim to make the following contributions:
\begin{enumerate}
    \item Use activation-aware SVD (actSVD) and difference-of-means (DoM) to test and gather empirical evidence
    about the relationship between safety and utility, specifically if they are orthogonal.
    \item Apply \citet{park2023LRH} causal separability framework to safety and utility.
    \item Provide a comparative analysis of actSVD and DoM as methods for identifying 
    and extracting directions in the activation space. 
\end{enumerate}